\documentclass[12pt,a4paper]{article}
\usepackage[T1]{fontenc}
\usepackage[utf8]{inputenc}
\usepackage{amsmath,amssymb}
\usepackage{booktabs}
\usepackage{tabularx}
\usepackage{geometry}
\geometry{top=25mm, bottom=25mm, left=28mm, right=28mm}
\usepackage{setspace}
\onehalfspacing
\usepackage[colorlinks=true,linkcolor=blue!60!black,citecolor=blue!60!black]{hyperref}

\begin{document}

% ============================================================
% INSERT POINT: This section follows Section 3.17 (Summary) in
% Chapter 3: "Application to Hybrid Intrusion Detection on
% Network-Traffic Graphs and Experimental Validation"
% ============================================================

\section{Production Deployment and Platform Validation}
\label{sec:platform-validation}

The preceding sections of this chapter established the theoretical
foundations and empirical performance of seven mathematical methods
(M1--M7) for hybrid intrusion detection on network-traffic graphs.
Validation on benchmark datasets---including CICIDS-2017, NSL-KDD,
UNSW-NB15, and CIC-Bell-DNS-2021---demonstrated that spectral graph
wavelets (M1), persistent homology (M2), Finsler geometry (M3),
sheaf-theoretic consistency (M4), category-theoretic composition (M5),
information-geometric detection (M6), and tropical semiring
aggregation (M7) each contribute complementary detection capabilities
with quantifiable improvements over conventional baselines [143].
However, benchmark evaluation alone is insufficient to establish
operational readiness. This section extends validation to a
production-grade platform, RobustIDPS.ai (DOI:
10.5281/zenodo.19129512) [144], which integrates all seven methods
alongside additional innovations required for deployment within
operational Security Operations Centres (SOCs). The platform serves
both as a validation instrument for the dissertation's theoretical
contributions and as an independent research artefact demonstrating
the viability of mathematically grounded intrusion detection at
production scale.

% ────────────────────────────────────────────────────────────
\subsection{Platform Architecture and Scope}
\label{subsec:platform-architecture}

RobustIDPS.ai is a full-stack, cloud-native intrusion detection and
response platform designed to unify the research methods presented in
this dissertation with the operational requirements of contemporary
SOC workflows. The platform encompasses \textbf{12+ neural network
models} maintained in a unified model registry, exposed through
\textbf{46 interactive pages} organised across \textbf{9 functional
groups}, as summarised in Table~\ref{tab:functional-groups}.

The system architecture follows a layered microservice design.  The
backend is implemented in \textbf{FastAPI} (Python 3.11+), providing
asynchronous REST endpoints and WebSocket channels for real-time
streaming.  The frontend is a \textbf{React 18} single-page
application with responsive design and Progressive Web App (PWA)
manifest, enabling cross-platform accessibility on desktop, tablet,
and mobile devices without requiring native application installation
[145].  Persistent state is managed by \textbf{PostgreSQL 16} with
connection pooling, while \textbf{Redis 7} serves as both a caching
layer and a pub/sub broker for inter-service communication.  The
entire stack is containerised via \textbf{Docker Compose}, with
separate service definitions for the API server, the frontend build,
the database, the cache, and background worker processes.  Deployment
is reproducible through a single \texttt{docker compose up} invocation
[146].

A distinguishing architectural feature is the \textbf{Multi-LLM SOC
Copilot}, which supports simultaneous integration with Claude
(Anthropic), GPT-4o (OpenAI), Gemini (Google), and DeepSeek via a
provider-agnostic adapter layer.  This design permits SOC analysts to
select the most appropriate large language model for a given
investigative task---such as alert summarisation, natural-language
threat hunting, or incident report generation---while maintaining a
consistent conversational interface [147].

The \textbf{Live Monitor} page implements WebSocket-based streaming
with a configurable retention window of up to 2{,}000 flows, allowing
analysts to observe classification decisions in real time as network
traffic is ingested.  Each flow is processed through the selected
model pipeline and annotated with predicted class, confidence score,
and MITRE ATT\&CK technique mapping before being pushed to the
client.  A companion \textbf{File Replay} mode enables deterministic
re-ingestion of stored PCAP captures for reproducible validation
experiments.

\begin{table}[htbp]
\centering
\caption{RobustIDPS.ai functional groups and page inventory.}
\label{tab:functional-groups}
\small
\begin{tabularx}{\textwidth}{l c X}
\toprule
\textbf{Functional Group} & \textbf{Pages} & \textbf{Key Capabilities} \\
\midrule
AI Command Centre      & 5 & Dashboard, Upload \& Analyse, Analytics, Live Monitor, Alert Causality \\
AI Active Defence      & 5 & Red Team Arena, Threat Response, SOC Copilot, Federated Learning, Attack Chain Predictor \\
SOC Intelligence       & 6 & Auto-Investigation, Threat Hunt, Incident Reports, Threat Intel, Rule Generator, CVE Mapper \\
LLM Attack Surfaces    & 5 & Prompt Injection, Jailbreak Taxonomy, RAG Poisoning, Multi-Agent Chain, Data Poisoning \\
MLSecOps Standards     & 3 & MITRE ATT\&CK Mapper, Alert Triage, Compliance Hub (OWASP + NIST) \\
AI Security \& Gov     & 4 & PQ Cryptography, Zero-Trust, Supply Chain, Research Hub \\
AI Novel Methods       & 5 & Continual Learning, RL Response Agent, Adversarial Eval, Explainability, Autoencoder Detector \\
AI Data \& Models      & 3 & Datasets, Models (12+), Ablation Studio \\
Industry \& Research   & 4 & Lab Partnerships, Postdoc Portal, Interview Prep, Benchmarks \\
\midrule
\textbf{Total}         & \textbf{40+} & \\
\bottomrule
\end{tabularx}
\end{table}

% ────────────────────────────────────────────────────────────
\subsection{Key Innovations Beyond the Core Methods}
\label{subsec:key-innovations}

While the seven mathematical methods (M1--M7) constitute the
theoretical core of this dissertation, their integration into a
production SOC platform necessitated six additional innovation
categories, each representing a substantive engineering and research
contribution.

\paragraph{1. LLM Security Testing Suite.}
The growing deployment of large language models within security
pipelines introduces novel attack surfaces that traditional IDS
architectures do not address [148].  RobustIDPS.ai includes four
dedicated modules for LLM security evaluation.  The \emph{Prompt
Injection} module implements a live \texttt{DefensePipeline} that
processes user inputs through three sequential stages: (i)~regex-based
pattern matching against 8~injection categories (direct override,
context manipulation, role hijacking, encoding attacks, nested
injection, multi-turn exploitation, system prompt extraction, and
output formatting attacks), (ii)~boundary enforcement that validates
input--output scope constraints, and (iii)~output filtering that
sanitises model responses before delivery to the analyst.  The
\emph{Jailbreak Taxonomy} module catalogues 8~jailbreak technique
families---including DAN-style prompts, character roleplay, hypothetical
framing, token smuggling, multi-language bypass, instruction hierarchy
exploitation, context window overflow, and reward hacking---and
provides live testing interfaces against real LLM endpoints to
measure defence efficacy [149].  The \emph{RAG Poisoning} module
demonstrates attacks against Retrieval-Augmented Generation pipelines
through document provenance tracking, Jaccard similarity-based
retrieval scoring, and poisoning detection heuristics that flag
documents with anomalous similarity distributions.  Finally, the
\emph{Multi-Agent Chain} module models a 4-agent topology (Planner,
Executor, Validator, Reporter) with inter-agent trust scoring,
enabling analysts to observe how adversarial manipulation of one agent
can cascade through the chain [150].

\paragraph{2. SOC Intelligence Suite.}
Drawing inspiration from commercial platforms such as Microsoft
Security Copilot, SentinelOne Purple AI, and Google Sec-Gemini
v1~[151], the SOC Intelligence suite provides six analyst-facing
tools.  \emph{Auto-Investigation} implements a 4-phase agentic
pipeline: (a)~alert enrichment via contextual feature extraction,
(b)~hypothesis generation using LLM-assisted reasoning,
(c)~evidence correlation across temporal and topological dimensions,
and (d)~verdict synthesis with confidence calibration.
\emph{Natural-Language Threat Hunt} exposes a query parser that
translates analyst questions into structured filters over attack type,
confidence threshold, port range, and severity level, thereby
lowering the barrier to proactive hunting for analysts without
specialised query-language training.  \emph{Incident Reports} are
auto-generated with MITRE ATT\&CK coverage annotations, executive
summaries, and technical appendices.  \emph{Threat Intel} provides IP
reputation scoring with geolocation enrichment.  The \emph{Rule
Generator} offers 13~Suricata/Snort rule templates covering common
attack signatures, enabling one-click rule synthesis from detected
anomalies.  The \emph{CVE Mapper} links 6~attack-class-to-CVE
mappings with associated CVSS v3.1 base scores, providing analysts
with vulnerability context for observed threats [152].

\paragraph{3. Multi-Agent PQC-IDS.}
A novel contribution extending the PQ-IDPS work (Branch~3 of the
dissertation) is the \emph{Multi-Agent PQC-IDS}, a cooperative
4-agent model comprising 70{,}598 trainable parameters [153].  The
\emph{Traffic Analyst} agent maps 83~input features to 34~attack
classes using a two-layer feed-forward network with batch
normalisation.  The \emph{PQC Specialist} agent focuses on
post-quantum cryptographic traffic, classifying 83~features into
14~PQ algorithm categories.  The \emph{Anomaly Detector} agent
employs an autoencoder with architecture $83 \to 16 \to 83$,
flagging flows whose reconstruction error exceeds a learned
threshold.  The \emph{Coordinator} agent performs attention-based
fusion over the three specialist outputs, producing a unified
detection verdict.  The model is publicly available at
\texttt{github.com/rogerpanel/Multi-Agent-PQC-models}, and the
training dataset is archived at DOI:
10.34740/kaggle/dsv/15424420~[154], comprising subsets from
CESNET-TLS22 and ArielCyber PQClass.

\paragraph{4. Continual Learning and Autonomous Response.}
To address concept drift in production traffic streams, the platform
implements Elastic Weight Consolidation (EWC) with experience replay
(buffer size $M = 5{,}000$ flows) [155].  The autonomous response
module uses Constrained Policy Optimisation (CPO) with 5~graduated
response actions ordered by severity: \textsc{Monitor} $\to$
\textsc{Log-Enrich} $\to$ \textsc{Rate-Limit} $\to$
\textsc{Isolate} $\to$ \textsc{Quarantine}.  A unified Fisher
Information matrix with blending coefficient $\beta = 0.7$ governs
the trade-off between plasticity and stability, while a bounded
autonomy configuration ensures that high-severity actions
(\textsc{Isolate}, \textsc{Quarantine}) require human-in-the-loop
confirmation before execution [156].

\paragraph{5. MLSecOps Standards Compliance.}
The platform provides structured compliance tooling aligned with three
major frameworks.  The \emph{MITRE ATT\&CK Mapper} associates each of
the 30~attack classes detected by the model ensemble with
corresponding ATT\&CK techniques and sub-techniques, enabling
threat-intelligence interoperability.  An \emph{OWASP LLM Top~10
Scorecard} evaluates the platform's own LLM integrations against each
risk category (LLM01--LLM10), with pass/fail indicators and
remediation guidance.  A \emph{NIST AI RMF Dashboard} tracks
alignment with the AI Risk Management Framework across its four
functions (Govern, Map, Measure, Manage).  The \emph{Alert Triage}
classifier applies a severity-weighted prioritisation algorithm that
combines model confidence, attack class severity, and contextual
indicators to rank alerts for analyst review [157].

\paragraph{6. Security Hardening.}
Production deployment demands defence-in-depth beyond the detection
algorithms themselves.  All user-facing content is sanitised through
\textbf{DOMPurify} to prevent cross-site scripting (XSS) attacks
against the analyst interface.  WebSocket connections require
\textbf{JWT token validation} at the handshake phase, preventing
unauthorised subscription to live traffic streams.  A \textbf{per-job
authorisation} model ensures that background tasks (model inference,
report generation) execute only under the identity of the requesting
user.  \textbf{Security headers middleware} enforces
Content-Security-Policy, X-Frame-Options, and
Strict-Transport-Security across all HTTP responses [158].

% ────────────────────────────────────────────────────────────
\subsection{Validation Datasets and Methodology}
\label{subsec:validation-datasets}

Validation of the integrated platform requires purpose-built datasets
that exercise the full breadth of detection capabilities.  Three
crafted PCAP captures were generated for this purpose, each designed
to test a distinct operational scenario:

\begin{enumerate}
    \item \textbf{validation\_benchmark.pcap} (323.5\,MB): Contains
    60{,}000 synthetic flows spanning 34~attack classes, with a
    60/40 split between labelled and unlabelled flows.  The labelled
    subset provides ground-truth annotations for supervised accuracy
    measurement, while the unlabelled subset tests the platform's
    unsupervised anomaly detection capabilities (M2, M6) and
    semi-supervised graph methods (M1, M4).

    \item \textbf{pqc\_validation\_benchmark.pcap}: A PQC-focused
    capture containing 40{,}000 flows generated from 10~post-quantum
    cryptographic algorithm implementations (Kyber-512, Kyber-768,
    Kyber-1024, Dilithium-2, Dilithium-3, Dilithium-5, SPHINCS+-128f,
    Falcon-512, Falcon-1024, Classic McEliece) and 4~PQ-specific
    attack variants (downgrade, oracle, hybrid-mode exploitation, and
    side-channel simulation).

    \item \textbf{combined\_validation\_benchmark.pcap}: An integrated
    capture of 50{,}000 flows with a deliberately heterogeneous
    composition: 40\% conventional IDS traffic, 30\% PQC-specific
    traffic, and 30\% adversarially perturbed flows designed to test
    robustness under evasion conditions.
\end{enumerate}

These crafted datasets are complemented by external PQC datasets
sourced from Kaggle (DOI: 10.34740/kaggle/dsv/15424420) [154],
including subsets from CESNET-TLS22 (Czech national research network
TLS captures), ArielCyber PQClass (post-quantum traffic
classification), PQS TLS Measurements (PQ algorithm handshake
timings), and PQ IoT Impact (resource-constrained device
benchmarks).  Together, these sources provide coverage across both
conventional and post-quantum threat landscapes.

The validation methodology follows a four-stage pipeline: (i)~upload
of PCAP files via the Live Monitor File Replay interface, which
deterministically re-ingests packets at configurable rates;
(ii)~multi-model detection, where each flow is classified by all
applicable models in the registry and results are aggregated through
the tropical semiring fusion layer (M7); (iii)~cross-page analysis,
in which detection outputs are correlated across the Analytics,
Alert Causality, and Auto-Investigation pages to assess end-to-end
SOC workflow coverage; and (iv)~comparison against ground-truth CSV
annotations, computing standard metrics (accuracy, precision, recall,
F1-score, and area under the ROC curve) per attack class and per
model [159].  This methodology ensures that validation captures not
only individual model performance but also the integrative behaviour
of the platform as a cohesive analytical instrument.

% === END SEGMENT 1 ===

\end{document}
